\documentclass[10pt,a4paper]{article}
\usepackage{amsmath,amssymb,amsthm,graphicx}
\usepackage{hyperref}
\usepackage{enumitem}
\usepackage{mathtools}
\title{An oracle blocks advice-free heuristic derandomization, with a conditional polynomial-space implementation}
\author{Shuichi Hirahara, Zhenjian Lu, Hanlin Ren}
\date{doi:10.4230/LIPIcs.CCC.2023.6; CC BY 4.0}
\begin{document}
\maketitle
\noindent\textit{Editorial note.} Original complete official proofs of11 and12 retained as a single major oracle-barrier result, preserving all original hypotheses and literal mathematical expressions. Earlier complete construction10 is context for author references; no supporting lemma or direct corollary is separately counted. Files unique to this task added; frozenCC030 dependencies unchanged.\par\medskip
\noindent\textit{Allowed standard prerequisites (editorial).}\begin{itemize}
\item Approximate counting in the polynomial hierarchy, padding and search-to-decision reductions. Exact original Theorems23/24 retained, including author counting sketch; standard padding under the unchanged SAT assumption used as in author proof
\end{itemize}\medskip
\paragraph{Definitions and notations (original Sections 1.1, 2.1).}
For a complexity class $\mathfrak C$, a complexity-theoretic statement is $\mathfrak C$-relativizing if it is true relative to every oracle $\mathcal O\in\mathfrak C$.
\paragraph{Definition 20.} A language $L$ is in $\mathsf{EXPH}$ if there is a constant $k$, an exponential bound $e(n)=2^{\operatorname{poly}(n)}$, and a $\operatorname{poly}(e(n))$-time algorithm $M$ such that, for every input $x\in\{0,1\}^*$,
\[x\in L\Longleftrightarrow\exists y_1\forall y_2\exists y_3\forall y_4\ldots\exists y_k,\ M(x,y_1,\ldots,y_k)=1,\]
where $y_i\in\{0,1\}^{e(n)}$ for each $i\in[k]$. Equivalently, a language is in $\mathsf{EXPH}$ if it can be decided in exponential ($2^{\operatorname{poly}(n)}$) time with access to a $\mathsf{PH}$ oracle.
\paragraph{Pseudodeterministic algorithms.} A pseudodeterministic algorithm is a probabilistic algorithm that returns a fixed output with high probability on any given input. We will also consider pseudodeterministic oracle (or advice-taking) algorithms. Such an algorithm will maintain its pseudodeterministic behavior (i.e., outputting a fixed answer with high probability for any input) when a correct oracle (or advice) is given.
\paragraph{Dense property (original Section 1.3.2).} A property $Q\subseteq\{0,1\}^*$ is dense if for every input length $n\in\mathbb N$, we have that $|Q\cap\{0,1\}^n|\geq2^n/\operatorname{poly}(n)$.
\paragraph{Theorem 23 [68].} For every oracle $\mathcal O$, $\mathsf{BPP}^{\mathcal O}\subseteq(\Sigma^p_2)^{\mathcal O}$.

\paragraph{Theorem24 (Approximate Counting in Quasi-Linear-Time Hierarchy).} There is a $\Sigma_5\mathsf{TIME}[n\cdot\operatorname{polylog}(n)]$ machine $M$ such that on input $(\varphi,K)$, accepts if $\varphi$ has at least $K$ satisfying assignments, and rejects if $\varphi$ has at most $K/2$ satisfying assignments.
\begin{proof}[Proof Sketch]
Using universal hashing [69], we can show that there is a $\mathsf{BPTIME}[n\cdot\operatorname{polylog}(n)]^{\mathsf{SAT}}$ machine $M'$ that approximates the number of satisfying assignments of $\varphi'$. Let $H$ be a family of pairwise-independent hash functions whose output range is $\{0,1\}^{\lceil\log K\rceil}$, that is computable in near-linear time. (See, e.g., [72, Problem3.3,(2)] for an example based on Toeplitz matrices.) The machine $M'$ randomly samples $h\sim H$ and $r\sim\{0,1\}^{\lceil\log K\rceil}$, and uses the $\mathsf{SAT}$ oracle to decide whether there is a satisfying assignment $x$ of $\varphi'$ such that $h(x)=r$.
Finally, the lemma follows from the fact that for every oracle $\mathcal O$, $\mathsf{BPTIME}^{\mathcal O}[t]\subseteq\Sigma_3\mathsf{TIME}[t\cdot\operatorname{polylog}(t)]^{\mathcal O}$ [74].
\end{proof}

\paragraph{Theorem 10.} There is an oracle $\mathcal O\in\mathsf{EXPH}$ such that
\[\mathsf{RP}^{\mathcal O}\not\subseteq\mathrm{i.o.}\,\mathsf{DTIME}^{\mathcal O}[2^n]\qquad\text{and}\qquad\mathsf{UP}^{\mathcal O}\not\subseteq\mathsf{BPTIME}^{\mathcal O}[2^n].\]

\begin{proof}[Proof of Theorem 10]
Our oracle world will consist of two oracles $P$ and $Q$ (it is easy to combine them into a single oracle). The oracle $P$ creates a hard problem in $\mathsf{RP}$ against deterministic algorithms, and the oracle $Q$ creates a hard problem in $\mathsf{UP}$ against probabilistic algorithms. In particular, every input to the oracle $P$ will be of the form $(x,r)$ where $|r|=100|x|$, and let
\[L_P:=\{x:\exists r\in\{0,1\}^{100|x|},P(x,r)=1\}.\]
We will guarantee that for every $x$, $\Pr_r[P(x,r)=1]$ is either $0$ or greater than $1/2$, thus $L_P\in\mathsf{RP}$. We also define
\[L_Q:=\{1^n:\exists r\in\{0,1\}^{100n},Q(r)=1\}.\]
We will guarantee that for every input length $100n$, $|Q\cap\{0,1\}^{100n}|\leq1$, thus $L_Q\in\mathsf{UP}$.
Let $M_1,M_2,\ldots$ be an enumeration of Turing machines running in deterministic $2^n$ time, and $N_1,N_2,\ldots$ be an enumeration of probabilistic Turing machines running in $2^n$ time (that does not necessarily satisfy the $\mathsf{BPP}$ promise).\footnote{In this paper, whenever we enumerate a list of Turing machines, we assume every machine occurs in the list infinitely many times.} We want every $M_i$ to fail to solve $L_P$ on all but finitely many input lengths, and every $N_i$ to fail to solve $L_Q$ on infinitely many input lengths. In particular, let $n_1$ be a large enough constant, and $n_i=2^{500n_{i-1}}$ for each $i\geq2$. Assuming that every machine appears in the sequence $\{N_i\}$ infinitely often, we want that for each $i\in\mathbb N$, $N_i$ fails to solve $L_Q$ on input $1^{n_i}$.
We also need the following terminologies. The $n$-th slice of $P$ is the set of inputs of the form $(x,r)$ where $|x|=n$ and $|r|=100n$. The $\ell$-th slice of $Q$ is simply the set of inputs with length $\ell$. During our construction, there will be some entries $P(x,r)$ and $Q(r)$ that are fixed, meaning that their values will never change in the subsequent construction. There will also be some inputs $x$ that are frozen, meaning that for every $r\in\{0,1\}^{100|x|}$, the values of $P(x,r)$ are fixed and will never change in the subsequent construction. If an entry $P(x,r)$ or $Q(r)$ is not fixed and, in the case of $P(x,r)$, $x$ is not frozen, then we say this entry is free.
Our construction proceeds in stages. During stage $n$, we will fix the entire $n$-th slices of $P$ and $Q$; we will possibly also fix or freeze some other entries. It will be guaranteed that before the $n$-th stage:
\begin{enumerate}
\item the number of frozen strings beyond the $n$-th slice (including the $n$-th slice) is at most $2^{n/2}$;
\item the number of fixed strings (except the frozen ones) beyond the $n$-th slice (including the $n$-th slice) is at most $2^{5n}$; and
\item each fixed or frozen entry beyond the $n$-th slice (including the $n$-th slice) is set to be $0$.
\end{enumerate}
\paragraph{Diagonalization against $M_i$.}
Let $x_1,x_2,\ldots,x_n$ be $n$ inputs of length $n$ that are not frozen. Since there are at most $2^{n/2}$ frozen inputs of length $n$, it is always possible to select $n$ unfrozen inputs. For each $i\in[n]$, we simulate the machine $M_i$ on input $x_i$. Whenever $M_i$ asks a query $P(x,r)$ or $Q(r)$, if this query is not free, then we return the corresponding value fixed before; if it is free, then we return $0$ as the answer and also fix this query to be $0$. Then $M_i$ will output an answer $\mathsf{ans}_i\in\{0,1\}$. Note that there are at most $2^{5n}+n\cdot2^n\ll2^{100n}/3$ entries in $P$ that are fixed, thus the vast majority of the entries $P(x_i,r)$ are free. We fix all these entries $P(x_i,r)$ to be $(1-\mathsf{ans}_i)$. It is easy to see that $M_i$ fails to compute $L_P$ on input $x_i$.
\paragraph{Diagonalization against $N_i$.}
If $n=10n_i$ for some $i\in\mathbb N$, then we also need to fix the $100n_i$-th slice of $Q$ so that $N_i$ fails to compute $L_Q$ on input $1^{n_i}$. (In stage $n$, this happens after the above diagonalization against $M_i$.) We simulate $N_i$ on input $1^{n_i}$. Note that $N_i$ is a probabilistic machine running in time $2^{n_i}$, thus it has $B:=2^{2^{n_i}}$ computational branches, where each branch has $1/B$ probability mass. On each branch, whenever $N_i$ asks a query $P(x,r)$ or $Q(r)$, if this query is not free, then we return the corresponding value fixed before; if it is free, then we return $0$ as the answer to the query. Note that we cannot fix the query as there are potentially as many as $B$ queries to fix. Instead, after simulating all the computational branches, we pick out the ``heavy'' queries that occur in a lot of branches and fix them:
\begin{itemize}
\item A string $x$ is heavy if at least $\frac{1}{10^4|x|^3}$ fraction of branches queried some free entry of the form $P(x,r)$.
\item An entry $Q(r)$ is heavy if at least $\frac{1}{10^4}$ fraction of branches queried $Q(r)$.
\end{itemize}
We freeze every heavy $x$ by setting every unfixed entry $P(x,r)$ to be $0$. We also fix every heavy $Q(r)$ to be $0$. Note that each computational branch of $N_i$ only makes at most $2^{n_i}$ queries. Therefore, for each input length $m$ where $n\leq m\leq2^{n_i}$, the number of heavy strings $x\in\{0,1\}^m$ is at most $10^4 2^{n_i}m^3$. The number of entries $Q(r)$ fixed before this stage is at most $2^{5n}$, and the number of heavy entries $Q(r)$ is at most $10^4 2^{n_i}$, therefore the total number of fixed entries in $Q$ after this stage is at most $2^{6n}\ll2^{100n_i}$.
After simulating every computational branch of $N_i$, let $p$ be the probability that $N_i$ outputs $1$ on input $1^n$. If $p<0.5$, then we pick a free $r\in\{0,1\}^{100n_i}$ and set $Q(r)=1$; since the number of fixed entries $Q(r)$ is $\ll2^{100n_i}$, such $r$ always exists. If $p\geq0.5$, then we do nothing.
Now we have that $N_i$ solves $L_Q$ on input $1^{n_i}$ with probability less than $2/3$ (actually, at most $1/2$). Is this always the case in the future? Consider how the future changes to the oracles might affect the behavior of $N_i$. Let $m$ be an input length where $n<m\leq2^n$. During stage $m$ (which has not happened yet), we will pick $m$ inputs $x_1,x_2,\ldots,x_m\in\{0,1\}^m$, then some entries of the form $P(x_i,r)$ will be set to be $1$. Also note that the probability mass of computational branches that each $x_i$ affects is at most $\frac{1}{10^4m^3}$. Therefore, the total probability mass of computational branches affected is at most
\[\sum_{m\in\mathbb N}\frac{m}{10^4m^3}<\frac{1}{10^3}.\]
Note that we also set at most one entry $Q(r)$ to be $1$, where $|r|=100n_i$. (The next time we diagonalize against $N_{i+1}$ happens in stage $n_{i+1}\gg2^n$.) Since this entry $Q(r)$ is not heavy, the probability mass of computational branches that it affects is at most $\frac{1}{10^4}$.
It follows that only an $\frac{11}{10^4}$ fraction of computational branches of $N_i$ will be affected in the future. That is, let $p'$ be the probability that $N_i$ outputs $1$ on input $1^{n_i}$ at the end of our construction (i.e., after stage $\gg2^n$), then $|p'-p|<\frac{11}{10^4}$. Therefore, it cannot be the case that
\[\Pr[N_i(1^{n_i})=L_Q(1^{n_i})]\geq2/3.\]
\paragraph{Clean-up.}
At the end of stage $n$, we fix every unfixed input $P(x,r)$ in the $n$-th slice to be $0$. If $n=10n_i$ then we also fix every unfixed input $Q(r)$ in the $10n$-th slice to be $0$. The number of frozen strings is at most $10^4 2^{n/10}n^3\leq2^{n/2}$ (note that at most one $N_i$ contributes to the frozen strings since the input lengths $\{n_i\}$ are very far apart). The number of fixed entries above the $(n+1)$-st slice is at most
\[2^{5n}+n\cdot2^n+\sum_{m=n}^{2^{n/10}}10^4 2^{n/10}m^3<2^{5(n+1)}.\]
\paragraph{The complexity of $\mathcal O$.}
We show that the above oracles $P$ and $Q$ are computable in $\mathsf{EXPH}$. It suffices to present a deterministic algorithm with access to a $\Sigma^p_3$ oracle that on input $1^n$, outputs the truth tables of the $n$-th slices of $P$ and $Q$ in $2^{O(n)}$ time.
We simulate each stage $n'\leq n$. Before stage $n'$, we maintain the set of fixed entries $P(x,r),Q(r)$, as well as the set of frozen strings $x$. There are at most $2^{O(n')}$ many such strings. In what follows, we use $\langle P,Q\rangle$ to denote (the length-$2^{O(n')}$ encoding of) the list of fixed entries and frozen strings.
It is easy to diagonalize against each $M_i$: we simply choose the inputs $x_1,x_2,\ldots,x_{n'}$, simulate each $M_i$ on input $x_i$, and fix the oracles accordingly. This takes $2^{O(n')}$ time.
To simulate $N_i$, we need to define the following language capturing the weight of queries. Consider simulating the probabilistic machine $N_i$ on the current oracle $\langle P,Q\rangle$. Whenever $N_i$ makes a query, if this query is not free, then we return the answer of the query; otherwise we return $0$. The following two claims are immediate corollaries of Theorem 23:
\paragraph{Claim 42.} There is a language $\mathsf{Heavy\text{-}String}(\langle P,Q\rangle,N_i,x)$ computable in $\Sigma^p_2$ such that:
\begin{itemize}
\item if at least $\frac{1.1}{10^4|x|^3}$ fraction of computational branches of $N_i$ queried some free entry of the form $P(x,-)$, then $\mathsf{Heavy\text{-}String}(\langle P,Q\rangle,N_i,x)=1$;
\item if at most $\frac{0.9}{10^4|x|^3}$ fraction of computational branches of $N_i$ queried some free entry of the form $P(x,-)$, then $\mathsf{Heavy\text{-}String}(\langle P,Q\rangle,N_i,x)=0$.
\end{itemize}
\paragraph{Claim 43.} There is a language $\mathsf{Heavy\text{-}Entry}(\langle P,Q\rangle,N_i,r)$ computable in $\Sigma^p_2$ such that:
\begin{itemize}
\item if at least $\frac{1.1}{10^4}$ fraction of branches of $N_i$ queried $Q(r)$, then $\mathsf{Heavy\text{-}Entry}(\langle P,Q\rangle,N_i,r)=1$;
\item if at most $\frac{0.9}{10^4}$ fraction of branches of $N_i$ queried $Q(r)$, then $\mathsf{Heavy\text{-}Entry}(\langle P,Q\rangle,N_i,r)=0$.
\end{itemize}
(Note that both $\mathsf{Heavy\text{-}String}$ and $\mathsf{Heavy\text{-}Entry}$ accepts inputs of length $2^{O(n')}$.)
With the aid of these two languages, it is possible to simulate stage $n'$ now. We say a string $x$ is heavy if $\mathsf{Heavy\text{-}String}(\langle P,Q\rangle,N_i,x)=1$, and an entry $Q(r)$ is heavy if $\mathsf{Heavy\text{-}Entry}(\langle P,Q\rangle,N_i,r)=1$. Since there are at most $2^{O(n')}$ heavy strings, we can find all of them in $2^{O(n')}$ time with access to an $\mathsf{NP}^{\mathsf{Heavy\text{-}String}}$ oracle. Similarly, since there are $2^{O(n')}$ heavy entries, we can find all of them in $2^{O(n')}$ time with access to an $\mathsf{NP}^{\mathsf{Heavy\text{-}Entry}}$ oracle. We freeze these heavy strings and fix these heavy entries accordingly.
Finally, we estimate the probability that $N_i$ outputs $1$ (given current versions of $P$ and $Q$) in $\Sigma^p_2$, within additive error $0.1$. If this estimation is less than $1/2$, then we pick a free $r\in\{0,1\}^{10n'}$ and set $Q(r)=1$; otherwise we do nothing.
At the end of stage $n'$, we fix every unfixed entry in the $n'$-th slice of $P(x,r)$ to be $0$. If $n=10n_i$ then we also fix every unfixed input $Q(r)$ in the $10n$-th slice to be $0$. This takes $2^{O(n')}$ time.
It is easy to see that simulating the $n$ stages takes $2^{O(n)}$ time with access to a $\Sigma^p_3$ oracle.
\end{proof}

\paragraph{Theorem11.} There is an oracle $\mathcal O$ such that
\[\mathsf{BPP}^{\mathcal O}\not\subseteq\mathrm{heur}\text{-}\mathsf{DTIME}^{\mathcal O}[2^n]\qquad\text{and}\qquad\mathsf{UP}^{\mathcal O}\not\subseteq\mathrm{i.o.}\,\mathsf{BPTIME}^{\mathcal O}[2^n].\]
\paragraph{Theorem12.} Suppose that $\mathsf{SAT}\in\mathsf{DTIME}[n\cdot\operatorname{polylog}(n)]\cap\mathsf{NC}$. Then there is an oracle $\mathcal O$ satisfying the conclusions of Theorem11 that can be computed in polynomial space.

\begin{proof}[Proof of Theorem11]
We use a similar construction as in Theorem10. The difference is that now we want a problem in $\mathsf{BPP}$ that is infinitely-often hard, and a problem in $\mathsf{UP}$ that is almost-everywhere hard, therefore we need to define the oracles differently. In particular, the oracle $P$ only receives one input $r$, and the oracle $Q$ receives two inputs $(x,r)$, where $|r|=10|x|$. The hard language in $\mathsf{BPP}$ is:
\[L_P=\{1^n:\Pr_{r\sim\{0,1\}^{n^4}}[P(r)=1]\geq1/2\}.\]
It is guaranteed that for every $n$, the fraction of length-$n^4$ strings that are in $P$ is either at most $1/3$ or at least $2/3$, therefore $L_P\in\mathsf{BPP}$. On the other hand, the hard language in $\mathsf{UP}$ is:
\[L_Q=\{x:\exists r\in\{0,1\}^{10|x|},Q(x,r)=1\}.\]
It is guaranteed that for every input $x$, there is at most one string $r\in\{0,1\}^{10|x|}$ such that $Q(x,r)=1$, therefore $L_Q\in\mathsf{UP}$. We will construct the oracles $P$ and $Q$ such that $L_P\notin\mathsf{DTIME}[2^n]$ and $L_Q\notin\mathrm{i.o.}\,\mathsf{BPTIME}[2^n]$. Note that $L_P$ is a unary language, therefore we also have $\mathsf{BPP}\not\subseteq\mathrm{heur}\text{-}\mathsf{DTIME}[2^n]$.
We use the same terminologies as before: $M_1,M_2,\ldots$ is an enumeration of deterministic Turing machines running in $2^n$ time, and $N_1,N_2,\ldots$ is an enumeration of probabilistic Turing machines running in $2^n$ time. The $\ell$-th slice of $P$ is the set of inputs with length $\ell$, and the $n$-th slice of $Q$ is the set of inputs of the form $(x,r)$ where $|x|=n$ and $|r|=10n$. An entry is fixed if its value will never change in the subsequent construction. We choose $n_1$ to be a large enough constant and define a sequence $\{n_i\}$ where $n_i=2^{10n_{i-1}}$ for each $i\geq2$. We want every $N_i$ to fail to solve $L_Q$ on all but finitely many input lengths, and every $M_i$ to fail to solve $L_P$ on input $1^{n_i}$.
Our construction proceeds in stages. We guarantee that:
\begin{itemize}
\item For every input length $n$ that is not in the sequence $\{n_i\}$, the $n$-th slice of $P$ is identically $0$.
\item For every input length $n$, there are at most $2n$ entries in the entire $n$-th slice of $Q$ that returns $1$.
\item Consider only the oracle $Q$. Before the $n$-th stage, the number of fixed entries beyond the $n$-th slice of $Q$ (including the $n$-th one) is at most $2^{6n}$, and all of them are fixed to be $0$.
\end{itemize}
We start from the $n^{\star}$-th stage for some large enough constant $n^{\star}$. For each $n\geq n^{\star}$, in the $n$-th stage, we diagonalize against the machines $N_i$ and $M_i$, and fix the $n$-th slice of $Q$ as well as the $n^4$-th slice of $P$ (when $n$ equals some $n_i$). Details follow.
\paragraph{Diagonalization against $N_i$.}
Let $n'$ be the smallest integer in the sequence $\{n_i\}$ such that $n'\geq n$, and $i'$ be the index such that $n'=n_{i'}$. The $(n')^4$-th slice of $P$ may contain either few or most strings, and we do not know which is the case yet. Therefore, we will diagonalize each $N_i$ twice, first assuming the $n'$-th slice of $P$ is nearly empty and then assuming the $n'$-th slice of $P$ is nearly full. For this reason, we will need $2n$ distinct strings $x_{i,b}\in\{0,1\}^n$, one for each pair of $i\in[n]$ and $b\in\{0,1\}$. Note that at the beginning of the $n$-th stage, there are at most $2^{6n}\ll2^{10n}$ strings fixed in the $n$-th slice of $Q$, therefore we can always choose $2n$ strings $x_{i,b}$ where no entry of the form $Q(x_{i,b},r)$ is fixed.
For each $i\in[n]$ and $b\in\{0,1\}$, we simulate $N_i(x_{i,b})$, assuming that most entries in the $n'$-th slice of $P$ returns $b$. Note that $N_i$ is a probabilistic machine running in time $2^n$, thus it has $B:=2^{2^n}$ computational branches, where each branch has $1/B$ probability mass. On each branch, whenever $N_i$ asks a query $P(r)$ or $Q(x,r)$, if this query is already fixed, then we return the corresponding value fixed before; otherwise:
\begin{itemize}
\item Suppose the query is $P(r)$. If $|r|$ is in the sequence $\{n_i^4\}$, then since $P(r)$ is not fixed and $|r|\leq2^n$, we know that $|r|=(n')^4$ and we return $b$ as the answer. Otherwise ($|r|\neq(n')^4$) we return $0$.
\item Suppose the query is $Q(x,r)$, then we return $0$ as the answer.
\end{itemize}
Again, we can only afford to fix the heavy entries, defined as follows. Let $K$ be the number of machine-input pairs that we still need to simulate in the future before fixing the $(n')^4$-th slice of $P$; note that this includes both the current $N_i(x_{i,b})$ and the final $M_{i'}(1^{n'})$. That is:
\[K=2(n-i)+(2-b)+\sum_{j=n+1}^{n'}(2j)+1.\]
\begin{itemize}
\item An entry $P(r)$ is heavy if $|r|=(n')^4$ and at least $2^{-(n'+4)K}$ fraction of branches queried $P(r)$.
\item An entry $Q(x,r)$ is heavy if at least $\frac{1}{10^4|x|^3}$ fraction of branches queried $Q(x,r)$.
\end{itemize}
We fix every heavy entry according to the way we answer this entry before. That is, every heavy entry $P(r)$ is fixed to be $b$ and every heavy entry $Q(x,r)$ is fixed to be $0$. The number of heavy entries in $P$ is at most $2^{(n'+4)K}\cdot2^n$; for each input length $m\geq n$, the number of heavy entries in the $m$-th slice of $Q$ is at most $2^n\cdot10^4\cdot m^3$.
Let $p$ be the probability (over the $B$ computational branches) that $N_i(x_{i,b})$ outputs $1$. If $p<1/2$ then we find a string $r$ such that $Q(x_{i,b},r)$ is unfixed yet and set $Q(x_{i,b},r)=1$, making $x_{i,b}\in L_Q$. Such a string $r$ exists since we have fixed strictly less than $2^{10n}$ strings in the $n$-th slice of $Q$. If $p\geq1/2$ then we do nothing.
Assuming that most entries in the $(n')^4$-th slice of $P$ indeed returns $b$, at this point the probability that $N_i$ outputs the correct answer on input $x_{i,b}$ is at most $1/2$. However, the oracles $P$ and $Q$ might change in the future, thus the behavior of $N_i$ might also change. What fraction of computational branches might be affected in the future?
\begin{itemize}
\item Consider the $K-1$ machine-input pairs that we will simulate in the future before the $(n')^4$-th slice of $P$ is completely fixed. The number of heavy entries in the $(n')^4$-th slice of $P$ that these machines fix is at most
\[\sum_{k=1}^{K-1}2^{(n'+4)k}\cdot2^{n'}\leq2\cdot2^{(n'+4)(K-1)}\cdot2^{n'}.\]
Therefore, the probability mass of computational branches of $N_i(x_{i,b})$ influenced by the future changes to the $(n')^4$-th slice of $P$ is at most
\[2\cdot2^{(n'+4)(K-1)}\cdot2^{n'}\cdot2^{-(n'+4)K}\leq1/8.\]
\item For each $m\leq2^n$, there are at most $2m$ entries in the $m$-th slice of $Q$ that we will fix to be $1$. The probability mass of computational branches of $N_i(x_{i,b})$ influenced by these entries is at most
\[\sum_{m=n}^{2^n}\frac{2m}{10^4m^3}<\frac{1}{10^3}.\]
\end{itemize}
It follows that the future modifications to the oracles $P$ and $Q$ will only affect the accept probability of $N_i(x_{i,b})$ by $\frac18+\frac{1}{10^3}<\frac16$. Therefore, it cannot be the case that
\[\Pr[N_i(x_{i,b})=L_Q(x_{i,b})]\geq2/3.\]
\paragraph{Diagonalization against $M_i$.}
If $n=n_i$ for some $i\in\mathbb N$, then we want that $M_i$ does not compute $L_P$ on input $1^n$ correctly. We simulate $M_i$ on the input $1^n$. When $M_i$ asks a query $P(r)$ or $Q(x,r)$, if this query is already fixed, then we return the corresponding value fixed before; otherwise we return $0$ and fix this queried entry.
After this simulation, we argue that only a small fraction of entries in the $n^4$-th slice of $P$ are fixed. First, the simulation of $M_i(1^n)$ fixes at most $2^n$ entries. Second, there are at most $K=\sum_{j=1}^n(2j)=O(n^2)$ machine-input pairs $N_i(x_{i,b})$ simulated so far, and the $k$-th such machine fixes at most $2^{(n+4)k}\cdot2^n$ entries. Therefore, the number of entries fixed by some previous $N_i(x_{i,b})$ is at most
\[\sum_{k=1}^K2^{(n+4)k}\cdot2^n\leq2^{O(n^3)}.\]
It follows that all but $2^{O(n^3)}$ entries in the $n^4$-th slice of $P$ are not fixed yet. Let $b\in\{0,1\}$ be the output bit of $M_i(1^n)$, then we set $L_P(1^n)=1-b$ by fixing every unfixed entry in the $n^4$-th slice of $P$ to be $1-b$. We have that on input $1^n$, $L_P$ satisfies the $\mathsf{BPP}$ promise, and $M_i$ fails to solve $L_P$.
\paragraph{Clear-up.}
At the end of stage $n$, we fix every unfixed input $Q(x,r)$ on the $n$-th slice to be $0$. It is easy to see that there are at most $2n$ entries in the $n$-th slice of $Q(x,r)$ that returns $1$.
During the $n$-th stage, the number of entries we fixed beyond the $(n+1)$-st slice of $Q$ is at most
\[2^n+\sum_{m=n+1}^{2^n}2^n\cdot10^4\cdot m^3\leq2^{5n}.\]
Thus the total number of entries we fixed beyond the $(n+1)$-st slice of $Q$ is at most $2^{6n}+2^{5n}<2^{6(n+1)}$.
\end{proof}

Unfortunately, we do not know how to compute the above oracles in $\mathsf{EXPH}$. The reason is that when we diagonalize $N_i$ on input length $n$, the number of heavy entries $P(r)$ we need to fix is $2^{(n'+4)K}\cdot2^n=2^{\operatorname{poly}(n')}$. Since $n'$ might be exponentially large compared to $n$, this upper bound might be doubly exponential. Nevertheless, we show that under the assumption that
\[\mathsf{SAT}\in\mathsf{DTIME}[n\cdot\operatorname{polylog}(n)]\cap\mathsf{NC},\tag{*}\]
the oracles $P,Q$ in Theorem11 can be computed in $\mathsf{PSPACE}$.
We sketch some intuition before presenting the full proof of Theorem12. The bottleneck of putting $\mathcal O$ into $\mathsf{EXPH}$ is that we do not have enough space to store all the heavy entries $P(r)$ fixed by a probabilistic machine $N_i$, as there might be doubly-exponentially many of them. Therefore, whenever we simulate a machine $N_i$ and it asks a query $P(r)$, we have to compute from scratch whether $P(r)$ was fixed by a previous machine. If we have an oracle $\mathsf{Heavy}$ that given $r$, decides whether $P(r)$ was already fixed, then we can simulate $N_i$ in exponential time with a constant number of alternations (using Theorem24). However, there are $\operatorname{poly}(n)$ machines $N_j$ simulated before $N_i$, and each time we invoke $\mathsf{Heavy}$, we need to simulate these machines to see if any of them has already fixed $P(r)$. It follows that the simulation of $N_i$ actually requires $\operatorname{poly}(n)$ alternations and $2^{\operatorname{poly}(n)}$ time. Still, under a strong enough assumption such as (*), we can simulate these $\operatorname{poly}(n)$ alternations in $2^{\operatorname{poly}(n)}$ time.
\begin{proof}[Proof of Theorem12]
Consider the construction in Theorem11, with the only difference that we do not calculate the precise probabilities; instead, we compute their approximations in $\mathsf{PH}$. Since $\mathsf{SAT}\in\mathsf{NC}$, we have $\mathsf{EXP}=\mathsf{PSPACE}$ by padding, thus it suffices to construct the oracles in $2^{\operatorname{poly}(n)}$ time.
The most involved part of this proof is to compute the following function $\mathsf{Heavy}(n,\langle P,Q\rangle,i,b,x_{i,b},r)$, which indicates whether $P(r)$ is a heavy query of the probabilistic machine $N_i$ on input $x_{i,b}$. Here, the input of $\mathsf{Heavy}$ consists of the description of oracles $P$ and $Q$, an integer $i\leq n$, an input $x_{i,b}\in\{0,1\}^n$, and a query $r\in\{0,1\}^{(n_{i'})^4}$ (where $n_{i'}$ is the smallest element in the sequence $\{n_i\}$ such that $n_{i'}\geq n$). We may assume $|r|\leq2^n$ as otherwise $N_i$ could never query $P(r)$. The oracle description $\langle P,Q\rangle$ will contain the following information:
\begin{itemize}
\item We record a table of all entries in $Q$ that are fixed; there are at most $\sum_{m\leq n}2m+2^{6n}\leq2^{O(n)}$ such entries. For each entry, we also record a timestamp indicating when this entry is fixed.
\item We record all entries of $P$ up to the $n_{i'-1}$-th slice (but we do not include any information of the $n_{i'}$-th slice of $P$); the description length of $P$ is also at most $2^{O(n)}$.
\end{itemize}
Note that the input length of $\mathsf{Heavy}$ is at most $2^{O(n)}$. Instead of requiring $\mathsf{Heavy}$ to decide exactly whether $P(r)$ is a heavy query, we only require a $2$-approximation: if $P(r)$ is queried by $N_i(x_{i,b})$ w.p. at least $2^{-(n'+4)K}$ then $\mathsf{Heavy}$ returns $1$, while if $P(r)$ is queried w.p. at most $2^{-(n'+4)K}/2$ then $\mathsf{Heavy}$ returns $0$.
Let $k$ be the number of machine-input pairs we have simulated before $N_i(x_{i,b})$ (inclusive). That is,
\[k:=\sum_{m<n}(2m)+2i+b+1\leq O(n^2).\]
We will show that:
\paragraph{Claim46.} If $\mathsf{SAT}\in\mathsf{DTIME}[n\cdot\operatorname{polylog}(n)]$, then we can compute $\mathsf{Heavy}$ in deterministic $O(2^{k^2})$ time.
\begin{proof}
Consider simulating (a random computational branch of) $N_i(x_{i,b})$ while answering the oracle queries of $N_i$ accordingly.\footnote{That is, whenever $N_i$ asks a query $P(r)$ or $Q(x,r)$ that is fixed, we answer accordingly. Whenever $N_i$ asks a query $P(r)$ that is not fixed, if $|r|=(n')^4$ then we return $b$, otherwise we return $0$. Whenever $N_i$ asks a query $Q(x,r)$ that is not fixed, we return $0$.} Note that we have not recorded the $n_{i'}$-th slice of $P$ in our description $\langle P,Q\rangle$. As a consequence, whenever $N_i(x_{i,b})$ asks a query $P(r_q)$ where $|r_q|=n_{i'}$, we need to recursively call $\mathsf{Heavy}$ to decide whether this query was already fixed and which value it was fixed to. We enumerate every $n_{i'-1}<\tilde n\leq n$, every machine $N_j$ ($j\leq\tilde n$), and every bit $b'$. Suppose that we feed the machine $N_j$ with the input $x'_{j,b'}\in\{0,1\}^{n_{i'}}$.\footnote{We assume that there is an easy and deterministic way of assigning the inputs $x_{j,b'}$ for each pair $(j,b')$.} Using the timestamps recorded in the table $Q$, we can recover the state of the oracle $Q$ before the simulation of $N_j(x'_{j,b'})$; we call this oracle $Q'$. (The portion of oracle $P$ up to the $n_{i'-1}$-th slice remains the same.) Then, we call $\mathsf{Heavy}(\tilde n,\langle P,Q'\rangle,j,b',x'_{j,b'},r_q)$ to see if this query is a heavy query fixed by $N'_j$. If it is, then we return $P(r_q):=b'$; otherwise we search through the next machine. If the query $P(r)$ was not fixed before and was asked in this simulation. then we return $1$; otherwise we return $0$.
The above argument implies that $\mathsf{Heavy}$ can be computed recursively in the following sense. There is a machine $V_{\mathsf{Heavy}}$ that gets $\mathsf{input}:=(n,\langle P,Q\rangle,i,b,x_{i,b},r)$ and some randomness $z\in\{0,1\}^{2^{O(n)}}$ (denoting a random computational branch of $N_i(x_{i,b})$), runs in $2^{O(n)}$ time with $\operatorname{poly}(n)$ invocations of the $\mathsf{Heavy}$ oracle, and outputs $0$ or $1$. Let $\varepsilon:=2^{-(n'+4)/K}$, we want that $\mathsf{Heavy}(\mathsf{input})=1$ if $\Pr_z[V_{\mathsf{Heavy}}(\mathsf{input};z)=1]\geq\varepsilon$, and $\mathsf{Heavy}(\mathsf{input})=0$ if $\Pr_z[V_{\mathsf{Heavy}}(\mathsf{input};z)=1]\leq\varepsilon/2$.
We can use induction to show that $\mathsf{Heavy}$ can be computed in (deterministic) $O(2^{k^2})$ time. In particular, by the induction hypothesis, each call of the $\mathsf{Heavy}$ oracle made by $V_{\mathsf{Heavy}}$ on $\mathsf{input}$ can be computed in $O(2^{(k-1)^2})$ time. Therefore, $V_{\mathsf{Heavy}}$ runs in at most $2^{(k-1)^2}\cdot\operatorname{poly}(k)$ time. By Theorem24, there is a $\Sigma_5\mathsf{TIME}[2^{(k-1)^2}\cdot\operatorname{poly}(k)]$ time machine $\widehat{\mathsf{Heavy}}$ such that $\widehat{\mathsf{Heavy}}(\mathsf{input})=1$ if $\Pr_z[V_{\mathsf{Heavy}}(\mathsf{input};z)=1]\geq\varepsilon$ and $\widehat{\mathsf{Heavy}}(\mathsf{input})=0$ if $\Pr_z[V_{\mathsf{Heavy}}(\mathsf{input};z)=1]\leq\varepsilon/2$. By our hypothesis that $\mathsf{SAT}\in\mathsf{DTIME}[n\cdot\operatorname{polylog}(n)]$, there is a deterministic machine $\mathsf{Heavy}$ that runs in $2^{(k-1)^2}\cdot\operatorname{poly}(k)<2^{k^2}$ time and decides the same language as $\widehat{\mathsf{Heavy}}$.
\end{proof}
Now we show that given an integer $N$, it is possible to print the $N$-th slice of $P$ and $Q$ in deterministic $2^{\operatorname{poly}(N)}$ time. We simulate the stages $n$ for $n=1,2,\ldots$, where during each stage we need to diagonalize against probabilistic machines $N_1,N_2,\ldots$. If $n=n_i$ for some $i\in\mathbb N$, then we also need to diagonalize against $M_i$.
\begin{itemize}
\item For each $N_i$, we can list the set of its heavy queries of the form $Q(x,r)$ in deterministic $2^{\operatorname{poly}(n)}$ time with a $\mathsf{PH}^{\mathsf{Heavy}}$ oracle. Note that $\mathsf{Heavy}$ can be decided in time quasi-polynomial in its input length; also recall we assumed that $\mathsf{SAT}\in\mathsf P$. Thus we can enumerate the heavy queries of the form $Q(x,r)$ in deterministic $2^{\operatorname{poly}(n)}$ time (without additional oracles). We fix all these queries. (Note that we do not fix the heavy queries of the form $P(r)$; instead, we use the oracle $\mathsf{Heavy}$ to decide whether a query $P(r)$ is fixed.) Then we use the $\mathsf{PH}$ oracle to estimate the probability that $N_i(x_{i,b})$ outputs $1$. If the probability is small then we pick some $r$ such that $Q(x_{i,b},r)$ is unfixed, and fix this entry to $1$; otherwise we do nothing.
\item Before we diagonalize against $M_i$, we spend $2^{\operatorname{poly}(n)}$ time to retrieve the list of fixed entries in the $n^4$-th slice of $P$ from the oracle $\mathsf{Heavy}$. Then we simulate $M_i$, and for every unfixed query, we fix it to be $0$. Let $b\in\{0,1\}$ be the output of $M_i$, then we set every unfixed entry in the $n^4$-th slice of $P$ to be $1-b$.
\end{itemize}
It is easy to see that the above process runs in $2^{\operatorname{poly}(N)}$ time.
\end{proof}

\begin{thebibliography}{99}
\bibitem[31]{GIPS} Shafi Goldwasser, Russell Impagliazzo, Toniann Pitassi, and Rahul Santhanam. On the pseudo-deterministic query complexity of NP search problems. CCC2021, LIPIcs200,36:1--36:22. doi:10.4230/LIPIcs.CCC.2021.36.
\bibitem[33]{H18} Shuichi Hirahara. Non-black-box worst-case to average-case reductions within NP. FOCS2018,247--258. doi:10.1109/FOCS.2018.00032.
\bibitem[34]{H20} Shuichi Hirahara. Non-disjoint promise problems from meta-computational view of pseudorandom generator constructions. CCC2020,LIPIcs169,20:1--20:47. doi:10.4230/LIPIcs.CCC.2020.20.
\bibitem[68]{Sipser} Michael Sipser. A complexity theoretic approach to randomness. STOC1983,330--335. doi:10.1145/800061.808762.
\bibitem[13]{BFT} Harry Buhrman, Lance Fortnow, and Thomas Thierauf. Nonrelativizing separations. CCC1998,8--12. doi:10.1109/CCC.1998.694585.
\bibitem[69]{Stockmeyer} Larry J. Stockmeyer. The complexity of approximate counting (preliminary version). STOC1983,118--126. doi:10.1145/800061.808740.
\bibitem[72]{Vadhan} Salil P. Vadhan. Pseudorandomness. Foundations and Trends in Theoretical Computer Science,7(1--3):1--336,2012. doi:10.1561/0400000010.
\bibitem[74]{Viola} Emanuele Viola. On approximate majority and probabilistic time. CCC2007,155--168. doi:10.1109/CCC.2007.16.
\end{thebibliography}
\end{document}
